\documentclass[11pt,a4paper]{article}

% ─── Paquetes ───────────────────────────────────────────────
\usepackage[utf8]{inputenc}
\usepackage[T1]{fontenc}
\usepackage{lmodern}
\usepackage[margin=2.3cm]{geometry}
\usepackage{amsmath,amssymb}
\usepackage{graphicx}
\usepackage{xcolor}
\usepackage{hyperref}
\usepackage{booktabs}
\usepackage{array}
\usepackage{tabularx}
\usepackage{enumitem}
\usepackage{fancyhdr}
\usepackage{titlesec}
\usepackage{listings}
\usepackage{tcolorbox}
\usepackage{fontawesome5}
\usepackage{parskip}

% ─── Colores (paleta de tonos tierra de la aplicación) ──────
\definecolor{accent}{HTML}{2E7D32}       % green (app accent)
\definecolor{brown}{HTML}{5D4037}        % dark brown (toolbar)
\definecolor{brownmid}{HTML}{6B4C3A}     % body background
\definecolor{brownlight}{HTML}{BCAAA4}   % panel
\definecolor{cream}{HTML}{EFEBE9}        % surface
\definecolor{paper}{HTML}{FFF8F0}        % plot background
\definecolor{warnhl}{HTML}{9A6700}
\definecolor{redhl}{HTML}{8B2E2E}
\definecolor{border}{HTML}{BCAAA4}

% ─── Configuración de Hyperref ──────────────────────────────
\hypersetup{
  colorlinks=true,
  linkcolor=accent,
  urlcolor=brown,
  citecolor=accent,
  pdftitle={Conditional Dimensional Explorer - Manual de usuario},
  pdfsubject={Manual de usuario}
}

% ─── Encabezado / pie de página ─────────────────────────────
\pagestyle{fancy}
\fancyhf{}
\fancyhead[L]{\small\textcolor{brown}{\textbf{Conditional Dimensional Explorer}}}
\fancyhead[R]{\small Edición en tonos tierra}
\fancyfoot[C]{\small\thepage}
\renewcommand{\headrulewidth}{0.4pt}

% ─── Estilo de secciones ────────────────────────────────────
\titleformat{\section}
  {\Large\bfseries\color{brown}}{\thesection}{0.6em}{}
\titleformat{\subsection}
  {\large\bfseries\color{black!80}}{\thesubsection}{0.5em}{}
\titleformat{\subsubsection}
  {\normalsize\bfseries\color{black!70}}{\thesubsubsection}{0.4em}{}

% ─── Recuadros personalizados ──────────────────────────────
\newtcolorbox{notebox}[1][]{
  colback=cream, colframe=brown, boxrule=0.6pt,
  arc=2mm, left=4mm, right=4mm, top=2mm, bottom=2mm,
  fonttitle=\bfseries\small, title={#1}
}
\newtcolorbox{tipbox}[1][]{
  colback=accent!6, colframe=accent, boxrule=0.6pt,
  arc=2mm, left=4mm, right=4mm, top=2mm, bottom=2mm,
  fonttitle=\bfseries\small, title={#1}
}
\newtcolorbox{warnbox}[1][]{
  colback=orange!6, colframe=warnhl, boxrule=0.6pt,
  arc=2mm, left=4mm, right=4mm, top=2mm, bottom=2mm,
  fonttitle=\bfseries\small, title={#1}
}

% ─── Listados de código ─────────────────────────────────────
\lstset{
  basicstyle=\ttfamily\small,
  backgroundcolor=\color{cream},
  frame=single,
  rulecolor=\color{border},
  breaklines=true,
  columns=fullflexible,
  keepspaces=true,
}

\setlength{\parindent}{0pt}
\setlength{\parskip}{5pt}

\begin{document}

% ════════════════════════════════════════════════════════════
% TITLE PAGE
% ════════════════════════════════════════════════════════════
\begin{titlepage}
\centering
\vspace*{3cm}

{\Huge\bfseries\color{accent} Conditional Dimensional Explorer}\\[0.4cm]
{\LARGE\bfseries Manual de usuario}\\[1.5cm]

{\large Edición en tonos tierra}\\[0.3cm]
{\large Herramienta interactiva de visualización condicional de datos}\\[3cm]

\begin{tabular}{rl}
\textbf{Tipo de aplicación:} & Aplicación web de una sola página (HTML) \\
\textbf{Motor de representación gráfica:} & Plotly.js 3.1.0 \\
\textbf{Vistas:} & Histograma, dispersión (2D), dispersión 3D \\
\textbf{Condicionamiento:} & Mediante controles deslizantes y filtro por umbral de distancia \\
\textbf{Entrada de datos:} & Carga de CSV y generadores de datos sintéticos \\
\textbf{Exportaciones:} & CSV \\
\end{tabular}

\vfill

{\small Este documento explica cómo utilizar Conditional Dimensional Explorer: cómo cargar o generar datos, condicionar la nube de puntos, alternar entre vistas unidimensionales, bidimensionales y tridimensionales, y exportar los resultados.}

\end{titlepage}

% ════════════════════════════════════════════════════════════
\tableofcontents
\newpage

% ════════════════════════════════════════════════════════════
\section{Introducción}
% ════════════════════════════════════════════════════════════

La herramienta \textbf{Conditional Dimensional Explorer} permite visualizar de forma interactiva y condicional datos multivariantes desde un navegador. Transforma cada variable a una escala uniforme $[0,\,1]$ mediante la función de distribución acumulada empírica (ECDF) y permite explorar subconjuntos de datos seleccionados mediante variables de \emph{condicionamiento}. Los puntos más próximos al punto de condicionamiento elegido se resaltan, mientras que los demás se vuelven progresivamente transparentes.

La aplicación ofrece tres modos de visualización, según el número de variables que se elijan para representar:

\begin{itemize}[leftmargin=1.4em]
  \item \textbf{1 variable} --- un histograma de los valores originales de la variable seleccionada.
  \item \textbf{2 variables} --- un diagrama de dispersión en el cuadrado unidad.
  \item \textbf{3 variables} --- un diagrama de dispersión tridimensional interactivo.
\end{itemize}

\subsection{Características principales}

\begin{itemize}[leftmargin=1.4em]
  \item Cargar cualquier archivo CSV numérico (separado por comas, punto y coma o tabuladores).
  \item Generar dos conjuntos de datos sintéticos: una muestra estructurada fija de 15 variables y una muestra realista correlacionada de 10 variables.
  \item Condicionar por cualquier subconjunto de variables mediante controles deslizantes.
  \item Atenuación según la distancia: el control deslizante «umbral de distancia» determina la rapidez con la que se atenúa la nube de puntos.
  \item Guardar el conjunto de datos activo como CSV.
  \item Ver los datos en una tabla y examinar una matriz de correlaciones de Pearson.
  \item Modo de pantalla completa y botón de salida.
\end{itemize}

\begin{notebox}[No se requiere instalación]
La aplicación consta de un único archivo HTML. Solo necesita conexión a Internet durante la primera carga, porque obtiene Plotly.js de una CDN. Todos los cálculos se realizan localmente en el navegador.
\end{notebox}

% ════════════════════════════════════════════════════════════
\section{Cómo funciona la visualización}
% ════════════════════════════════════════════════════════════

\subsection{Normalización mediante la función de distribución acumulada empírica}

Cada variable se transforma a una escala uniforme en $[0,1]$ mediante su \textbf{función de distribución acumulada empírica} (ECDF):
\begin{equation}
  q_{ij} = \frac{\operatorname{rank}(x_{ij})}{n+1}
\end{equation}
donde $n$ es el número de filas y $\operatorname{rank}(x_{ij})$ es el rango del valor de la columna $j$ en la fila $i$ (el rango 1 corresponde al menor valor). De este modo, todos los ejes son comparables, independientemente de las unidades o de la escala.

\subsection{Condicionamiento y umbral de distancia}

Si selecciona $k$ variables de condicionamiento $\{c_1,\dots,c_k\}$, aparecerán también $k$ controles deslizantes en el panel derecho. Para cada punto se calcula una distancia en el espacio normalizado:

\begin{equation}
  d_i = \sqrt{\sum_{m=1}^{k} \left(q_{i c_m} - \theta_{c_m}\right)^2}
\end{equation}
donde $\theta_{c_m}$ es el valor del control deslizante de la variable de condicionamiento. Los puntos que cumplen $d_i < \tau$ (siendo $\tau$ el \emph{umbral de distancia}) se muestran con una opacidad

\begin{equation}
  \alpha_i = \max\left(0.2,\; 1 - \frac{d_i}{\tau}\right).
\end{equation}

Los puntos con $d_i \ge \tau$ se ocultan. Al reducir $\tau$, la atención se concentra en un entorno pequeño; al aumentarlo, se muestran más puntos.

\subsection{Los tres modos de visualización}

\begin{center}
\begin{tabularx}{\textwidth}{lX}
\toprule
\textbf{Variables representadas} & \textbf{Qué se muestra} \\
\midrule
1 variable & Un histograma de los valores originales (sin normalizar) de esa variable, filtrado según el condicionamiento actual. \\
2 variables & Un diagrama de dispersión 2D de sus coordenadas ECDF $[0,1]\times[0,1]$, con un marco verde. \\
3 variables & Un diagrama de dispersión 3D de sus coordenadas ECDF dentro de un cubo unidad. \\
\bottomrule
\end{tabularx}
\end{center}

\begin{tipbox}[En 2D y 3D, los ejes siempre son coordenadas ECDF]
En las vistas 2D y 3D, los ejes muestran valores \textbf{normalizados} entre 0 y 1, por lo que las posiciones de los puntos son comparables en todos los ejes. En la vista de histograma (1 variable), el eje x muestra los valores originales.
\end{tipbox}

% ════════════════════════════════════════════════════════════
\section{Primeros pasos}
% ════════════════════════════════════════════════════════════

\subsection{Abrir la aplicación}

Abra el archivo HTML en un navegador moderno (Chrome, Firefox, Edge o Safari). En la primera carga, la aplicación genera automáticamente la muestra estructurada fija de 15 variables, de modo que puede empezar a explorar de inmediato.

\subsection{Descripción de la interfaz}

\begin{center}
\begin{tabularx}{\textwidth}{lX}
\toprule
\textbf{Zona} & \textbf{Contenido} \\
\midrule
Barra de herramientas superior & Título, carga de archivos, tamaño de muestra, generadores de datos, selección de variables y botones para guardar, mostrar la tabla, ver correlaciones, activar pantalla completa y salir. \\
Gráfico central & La figura activa (histograma, dispersión o gráfico 3D). \\
Panel derecho & Umbral de distancia, contador de puntos y controles deslizantes de condicionamiento. \\
Ventanas modales & Selección de variables, tabla de datos y matriz de correlaciones. \\
\bottomrule
\end{tabularx}
\end{center}

\subsection{El contador de puntos}

Debajo del control deslizante del umbral de distancia aparece un indicador del tipo \texttt{2000 / 2000 pts}. El primer número indica cuántos puntos cumplen actualmente el umbral de distancia; el segundo es el número total de filas. Así puede comprobar inmediatamente lo restrictivo que es el condicionamiento.

% ════════════════════════════════════════════════════════════
\section{La barra de herramientas en detalle}
% ════════════════════════════════════════════════════════════

\begin{center}
\begin{tabularx}{\textwidth}{lX}
\toprule
\textbf{Botón} & \textbf{Acción} \\
\midrule
\faFolderOpen{} \textbf{Cargar CSV} & Abre un selector de archivos. Carga un CSV numérico y detecta automáticamente si está separado por comas, punto y coma o tabuladores. \\
\textbf{n =} (field) & Tamaño de muestra para los generadores sintéticos (200--20\,000). \\
\faBolt{} \textbf{Inicializar datos} & Genera el conjunto de datos estructurado fijo de 15 variables. \\
\faDna{} \textbf{Realista \textpm\ corr (10 vars)} & Genera una muestra correlacionada de 10 variables (Ingresos, Educación, Deuda, \dots). \\
\faSlidersH{} \textbf{Selección de variables} & Abre la ventana modal para elegir las variables representadas y las de condicionamiento. \\
\faSave{} \textbf{Guardar CSV} & Descarga el conjunto de datos actual. \\
\faClipboard{} \textbf{Mostrar tabla de datos} & Abre una ventana modal con las primeras 500 filas. \\
\faChartLine{} \textbf{Matriz de correlaciones} & Abre un mapa de calor de las correlaciones de Pearson. \\
\faExpand{} \textbf{Pantalla completa} & Activa o desactiva la pantalla completa. \\
\faPowerOff{} \textbf{Exit} & Intenta cerrar la pestaña (sujeto a las restricciones del navegador). \\
\bottomrule
\end{tabularx}
\end{center}

\begin{notebox}[El tamaño de muestra solo afecta a los generadores]
El campo \texttt{n =} determina cuántas filas generan los dos generadores sintéticos. No afecta a los archivos CSV que cargue; el campo se actualiza para reflejar el número de filas del archivo cargado.
\end{notebox}

% ════════════════════════════════════════════════════════════
\section{Carga de datos}
% ════════════════════════════════════════════════════════════

\subsection{Desde un archivo CSV}

Haga clic en \textbf{Cargar CSV} y seleccione un archivo. La aplicación:

\begin{enumerate}[leftmargin=1.4em]
  \item Detecta el separador (coma, punto y coma o tabulador) a partir de la primera línea no vacía.
  \item Analiza el archivo (admite campos entre comillas y saltos de línea entrecomillados).
  \item Conserva únicamente las filas cuyas celdas sean todas numéricas y estén completas.
  \item Si hay filas no válidas, muestra un cuadro de confirmación con los primeros números de fila afectados antes de cargar el subconjunto válido.
\end{enumerate}

La primera fila debe contener la cabecera (nombres de las variables); cada fila posterior debe tener exactamente tantos valores numéricos como nombres haya en la cabecera.

\subsubsection*{Coma decimal}

Cuando el archivo está separado por punto y coma o tabuladores, se admiten comas decimales (por ejemplo, \texttt{3,14}). En los archivos separados por comas, el separador decimal debe ser un punto.

\subsection{Sustituir el conjunto de datos activo}

Cada vez que cargue o genere un conjunto de datos nuevo, si ya hay uno activo la aplicación preguntará si desea guardarlo primero como CSV. Así se evita perder trabajo accidentalmente.

\begin{warnbox}[Sustitución irreversible]
Si decide no guardarlo, se descartarán el conjunto de datos activo, sus gráficos, sus ventanas modales y el estado de condicionamiento. Los datos nuevos se cargarán con un estado inicial: dos variables representadas, una variable de condicionamiento, umbral 0.80 y todos los controles deslizantes en 0.50.
\end{warnbox}

\subsection{Conjuntos de datos sintéticos}

Hay dos generadores integrados disponibles:

\begin{itemize}[leftmargin=1.4em]
  \item \textbf{Inicializar datos} --- una estructura fija de 15 variables con nombres de materias académicas (matemáticas, física, química, \dots, coordinación), generada mediante un modelo latente de cinco factores, con una matriz de cargas fija y una pequeña matriz de factores correlacionados.
  \item \textbf{Realista \textpm\ corr} --- 10 variables (ingresos, educación, deuda, ahorros, estrés, salud, puntuación social, tolerancia al riesgo, horas de trabajo y tiempo libre), generadas a partir de una matriz de correlación objetivo de dimensión $10\times10$, ajustada para que sea definida positiva mediante una pequeña regularización y muestreadas después mediante un factor de Cholesky.
\end{itemize}

Ambos generadores limitan todos los valores al intervalo $[0,100]$ y los redondean a una cifra decimal.

% ════════════════════════════════════════════════════════════
\section{Selección de variables}
% ════════════════════════════════════════════════════════════

Haga clic en \textbf{Selección de variables} para abrir la ventana modal. Contiene dos listas:

\begin{itemize}[leftmargin=1.4em]
  \item \textbf{Variables representadas} --- una, dos o tres. Son las variables que se representan en el gráfico.
  \item \textbf{Variables de condicionamiento} --- cualquier número. Cada una tendrá su propio control deslizante.
\end{itemize}

\subsection{Reglas}

\begin{itemize}[leftmargin=1.4em]
  \item Debe seleccionar al menos una variable representada.
  \item Puede seleccionar como máximo tres variables representadas.
  \item Las variables de condicionamiento no pueden coincidir con las variables representadas; la aplicación desactiva esas opciones en la lista de condicionamiento.
  \item Puede elegir \textbf{Ninguna (sin condicionamiento)} para eliminar todos los controles deslizantes de condicionamiento.
\end{itemize}

Click \textbf{Apply \& close} to accept. The app rebuilds the sliders and the plot immediately, and then requests fullscreen (a convenience feature that some browsers may block if no user gesture is detected).

\begin{tipbox}[Flujos de trabajo habituales]
\textbf{One-dimensional}: represente una variable y condicione por otras 2--3 para observar cómo cambia su distribución. \\
\textbf{Two-dimensional}: represente dos variables y condicione por otras 1--2 para observar la forma de la nube en secciones próximas. \\
\textbf{Three-dimensional}: represente tres variables, sin condicionamiento, para examinar el cubo completo.
\end{tipbox}

% ════════════════════════════════════════════════════════════
\section{Condicionamiento y umbral}
% ════════════════════════════════════════════════════════════

Cuando haya al menos una variable de condicionamiento, su control deslizante aparecerá en el panel derecho. Cada control varía entre 0 y 1 (unidades ECDF).

\subsection{El umbral de distancia}

El control deslizante superior del panel derecho, denominado \emph{umbral de distancia}, varía entre 0 y 3. Su valor $\tau$ determina:

\begin{itemize}[leftmargin=1.4em]
  \item \textbf{Puntos ocultos}: todos los puntos que cumplan $d_i \ge \tau$.
  \item \textbf{Opacidad}: $\alpha_i = \max(0.2,\, 1 - d_i/\tau)$ for the points kept.
\end{itemize}

Por tanto, un valor \textbf{menor} de $\tau$ produce un enfoque más estricto (menos puntos, más destacados y próximos al punto de condicionamiento), mientras que un valor \textbf{mayor} produce un enfoque más amplio (más puntos y una opacidad más uniforme).

\subsection{El contador de puntos}

El indicador situado debajo del umbral muestra \texttt{kept / total pts} (puntos conservados / puntos totales). Utilícelo para cuantificar lo restrictiva que es la configuración actual.

\begin{tipbox}[Cómo obtener un enfoque útil]
Empiece con el umbral en 0.80. Si el indicador muestra demasiados puntos (por ejemplo, más de la mitad del conjunto de datos), reduzca el umbral; si muestra muy pocos, auméntelo. Al mover un control de condicionamiento, el indicador se actualiza inmediatamente.
\end{tipbox}

% ════════════════════════════════════════════════════════════
\section{Los tres modos de visualización}
% ════════════════════════════════════════════════════════════

\subsection{Histograma (1 variable representada)}

El histograma cuenta los valores originales, sin normalizar, de esa variable entre los puntos que superan el filtro de condicionamiento. El eje x muestra los valores originales, no las coordenadas ECDF. Es la forma más rápida de ver cómo responde la distribución de una variable al condicionamiento.

\subsection{Diagrama de dispersión (2 variables representadas)}

Los puntos se muestran dentro del cuadrado unidad $[0,1]\times[0,1]$ en coordenadas ECDF, con un marco verde que indica los límites. La transparencia de cada punto refleja su proximidad al punto de condicionamiento; los límites de los ejes son siempre $[0,1]$.

\subsection{Diagrama de dispersión 3D (3 variables representadas)}

Un diagrama de dispersión 3D interactivo de las coordenadas ECDF dentro del cubo unidad. Los controles de cámara de Plotly están habilitados (arrastre para girar y use la rueda para acercar o alejar). La posición de la cámara se conserva al mover los controles deslizantes, de modo que puede centrarse en la región de interés sin que la vista cambie bruscamente en cada actualización.

\begin{notebox}[La apariencia está fijada]
En las vistas 2D y 3D, los ejes siempre están fijados en $[0,1]$ para que las comparaciones entre secciones sigan siendo significativas. Para restablecer la vista, utilice el botón «Home» de la barra de herramientas de Plotly.
\end{notebox}

% ════════════════════════════════════════════════════════════
\section{La tabla de datos y la matriz de correlaciones}
% ════════════════════════════════════════════════════════════

\subsection{Tabla de datos}

Haga clic en \textbf{Mostrar tabla de datos} para ver las primeras 500 filas. Las columnas conservan exactamente los nombres de la cabecera y los valores numéricos se redondean a tres decimales para facilitar la lectura. Si el conjunto contiene más de 500 filas, una fila al pie indica cuántas se han ocultado.

\subsection{Matriz de correlaciones}

Haga clic en \textbf{Matriz de correlaciones} para calcular la correlación de Pearson de cada par de variables y mostrarla en un mapa de calor. La escala va de $[-1,1]$ (de azul a rojo) y cada celda incluye el valor numérico. Esto permite comprobar si una variable de condicionamiento está correlacionada o no con una variable representada antes de seleccionarla.

% ════════════════════════════════════════════════════════════
\section{Guardar el trabajo}
% ════════════════════════════════════════════════════════════

Haga clic en \textbf{Guardar CSV} para descargar el conjunto de datos actual en un archivo separado por comas. Se conservan todos los nombres de las variables y cada valor se escribe sin modificaciones, de modo que al volver a importarlo se obtienen los mismos datos. Se añade una marca de orden de bytes UTF-8 al principio para que Excel reconozca la codificación.

\begin{warnbox}[Guardar antes de generar]
Al cargar un CSV o pulsar un generador, se le preguntará si desea guardar los datos activos, si los hay. Acepte siempre el aviso si quiere conservar su trabajo.
\end{warnbox}

% ════════════════════════════════════════════════════════════
\section{Pantalla completa y salida}
% ════════════════════════════════════════════════════════════

El botón \textbf{Fullscreen} activa o desactiva el modo de pantalla completa para toda la página. También se activa automáticamente al aplicar una nueva selección de variables.

El botón \textbf{Exit} solicita confirmación e intenta cerrar la pestaña. Por motivos de seguridad del navegador, una página solo puede cerrarse a sí misma si fue abierta mediante un script. Si no puede cerrarse, aparece un aviso que recuerda que debe cerrar la pestaña manualmente.

% ════════════════════════════════════════════════════════════
\section{Solución de problemas}
% ════════════════════════════════════════════════════════════

\begin{description}[leftmargin=0em, style=nextline]
  \item[``Plotly library could not be loaded'']
  La página no puede acceder a la CDN de Plotly. Conéctese a Internet y vuelva a cargar la página. Sin Plotly, la aplicación no puede generar gráficos.

  \item[``El archivo debe contener al menos dos variables'']
  El CSV solo tiene una columna. Asegúrese de que la fila de cabecera contiene al menos dos nombres.

  \item[``No valid numeric rows found'']
  Una o varias filas no contienen exclusivamente valores numéricos. Si utiliza comas decimales, asegúrese de que el archivo esté separado por punto y coma o tabuladores, no por comas.

  \item[``Unbalanced quotes'']
  A un campo entrecomillado le falta la comilla de cierre. Corrija el CSV e inténtelo de nuevo.

  \item[Solo se ven unos pocos puntos]
  El umbral de distancia es demasiado pequeño. Auméntelo o mueva los controles deslizantes de condicionamiento hacia el centro (0.5).

  \item[La variable seleccionada no aparece como eje]
  Probablemente se encuentra en la lista de condicionamiento. La aplicación desactiva las opciones coincidentes en la ventana modal; vuelva a abrirla y ajuste la selección.

  \item[La pantalla completa no se abre automáticamente]
  Algunos navegadores bloquean la pantalla completa automática si no se activa mediante una acción directa del usuario. En ese caso, utilice el botón Pantalla completa.
\end{description}

% ════════════════════════════════════════════════════════════
\section{Referencia rápida}
% ════════════════════════════════════════════════════════════

\begin{center}
\begin{tabularx}{\textwidth}{lX}
\toprule
\textbf{Tarea} & \textbf{Acción} \\
\midrule
Cargar un CSV & Cargar CSV \textrightarrow{} seleccionar archivo. \\
Generar datos predeterminados & Inicializar datos. \\
Generar datos correlacionados & Realista \textpm\ corr (10 vars). \\
Cambiar las variables representadas & Selección de variables \textrightarrow{} pick 1, 2 or 3 \textrightarrow{} Apply. \\
Cambiar el condicionamiento & Same modal; pick conditioning variables. \\
Enfocar más & Reducir el control deslizante del umbral de distancia. \\
Examinar filas & Mostrar tabla de datos. \\
Examinar correlaciones & Matriz de correlaciones. \\
Exportar datos & Guardar CSV. \\
Pantalla completa & Botón Pantalla completa (o automáticamente después de Apply). \\
Salir & Botón Exit. \\
\bottomrule
\end{tabularx}
\end{center}

\begin{notebox}[Flujo de trabajo recomendado]
\begin{enumerate}
  \item Cargue o genere el conjunto de datos.
  \item Examine la matriz de correlaciones para comprender las relaciones.
  \item Abra Selección de variables y, para empezar, elija dos variables representadas y una variable de condicionamiento.
  \item Mueva el control deslizante de condicionamiento para explorar distintas secciones.
  \item Reduzca el umbral de distancia para centrarse en un entorno local.
  \item Cambie a tres variables representadas para ver el cubo completo.
  \item Guarde el CSV si ha modificado los datos.
\end{enumerate}
\end{notebox}

% ════════════════════════════════════════════════════════════
\section*{Notas de la versión}
\addcontentsline{toc}{section}{Notas de la versión}
% ════════════════════════════════════════════════════════════

Este manual corresponde a la versión guardada como \texttt{ConditionalExplorerFINALNEW.html}. Sus características distintivas respecto de versiones anteriores incluyen:

\begin{itemize}[leftmargin=1.4em]
  \item Sustitución completa al cargar un conjunto de datos nuevo (se borran el estado anterior, los gráficos y las ventanas modales).
  \item Sustitución segura del CSV con una pregunta opcional para guardar antes de reemplazar.
  \item Analizador CSV robusto, con detección automática del separador y validación de cada fila.
  \item Conservación de la posición de la cámara de Plotly al actualizar los controles deslizantes en la vista 3D.
  \item Opción «Ninguna (sin condicionamiento)» en la lista de condicionamiento.
  \item Activación automática de la pantalla completa al aplicar una nueva selección de variables.
  \item Funcionamiento del umbral de distancia: los puntos con $d_i \ge \tau$ se ocultan y la opacidad disminuye según $1 - d_i/\tau$.
\end{itemize}

\end{document}